\section{Introduction}
\label{sec:Int}
\par The introduction is intended to give a more detailed overview of the research from the abstract. Document the \textbf{theme}, the \textbf{aim} and \textbf{motivation}. By motivation we mean the justification to why this project/research is relevant/needed for the current year/period and for your specific area of studies. You can consider also including your \textbf{hypothesis} and \textbf{research questions} in this section. Conclude this section with an overview of what to expect in the next sections. It is good practice to cross-reference your own paper using the ref command for the matching section label, such as in Section~\ref{sec:Lit} a review of current research is found.
\par Due to shift towards “Industry 4.0”, companies were forced to seek automated solutions for hiring. With rise of Transformer-based models using Natural Language Inference (NLI), it allows machines to understand the intent behind a candidate's story. The aim of this research was to develop and validate an interview prototype utilizing Zero-Shot NLP to provide real-time scoring of candidates personality traits. Traditional interview are prone to bias and expenses to run an interview. With use of BART-MNLI, it prevents lacking qualitative depth of conversation while employing automated scoring based on required labels from open-ended answers. 